% Adición a las conclusiones de VII; las conclusiones anteriores se conservan.
\paragraph{Integration of the four interactions and total canonical closure.}
Gravitation and the strong, weak and electromagnetic interactions act on the same material realization of the APP--TRIT--TPK structure. Transport preserves the sheets, orientation, memory and incidence of the enriched state; its spin and gauge realizations compose on the same material fibres. The electromagnetic component belongs to the electroweak realization: it is not introduced as an additional interaction disconnected from it. The total action uses the coefficients and readers already constructed, including the Planck section, radial length and Barbero--Immirzi functional.

The content of this integration is not limited to collecting terms in an expression. Variation of the same action produces the total spin current and the joint metric source. Eliminating contorsion preserves the cross-terms between species; the Legendre transformation preserves material momenta and produces the total constraints. The proofs in sections~\ref{hmtcuatro:gea:seccion} and~\ref{hmtcuatro:accion-retroaccion-restricciones} establish backreaction, the first-class property and propagation of those constraints in the specified regular realization.

In particular, let $\Theta_{\rm tot}$ be the total canonical potential obtained in that transformation, $\mathcal Z=\{C_A=0\}$ the regular constraint surface, and $X_N=N^AX_{C_A}$, $X_M=M^BX_{C_B}$ two admissible characteristic deformations. The connection induced by the action has $H_N^{\rm tot,can}=-\Theta_{\rm tot}(X_N)$. Its curvature, calculated without omitting derivatives of the deformation coefficients, is
\begin{equation}
\begin{aligned}
\mathcal F^{\rm tot,can}_{N,M}
&=\delta_NH_M^{\rm tot,can}-\delta_MH_N^{\rm tot,can}
 -H_{[X_N,X_M]}^{\rm tot,can}
 +\frac{i}{\hbar}[H_N^{\rm tot,can},H_M^{\rm tot,can}]\\
&=-\mathrm d\Theta_{\rm tot}(X_N,X_M)
 =N^AM^B f_{AB}{}^{D}C_D.
\end{aligned}
\label{hmtcuatro:conclusion-curvatura-total}
\end{equation}
Therefore,
\[
 \left.\mathcal F^{\rm tot,can}_{N,M}\right|_{\mathcal Z}=0.
\]
The formula includes the tangential deformation and the gauge actions of the total bracket; on the quotient, it expresses the same equality modulo those vertical directions. The prequantum representation preserves this closure through the commutator identity proved in \eqref{hmtcuatro:representacion-precuantica}.

Thus, gravitational incorporation is part of the same action, its sources and its constraint algebra, not an equation juxtaposed after constructing the other interactions. The proved vanishing concerns the canonical characteristic curvature; it does not require spacetime curvature or gauge-field strengths to vanish. Its domain preserves the nondegenerate coframe, the regular chart and the boundary conditions used in the proof. Nor does it suppress the memory or global monodromy of the histories that gave rise to the realization.
